% Rozdział w polskiej monografii
% Sąsiedztwo Motzkina dla permutacyjnego problemu przepływowego
\documentclass[12pt,a4paper]{article}
\usepackage[utf8]{inputenc}
\usepackage[T1]{fontenc}
\usepackage[polish]{babel}
\usepackage{amsmath,amssymb,amsthm}
\usepackage{booktabs}
\usepackage{graphicx}
\usepackage{tikz}
\usepackage{hyperref}
\usepackage{algorithm}
\usepackage{algpseudocode}

\newtheorem{definicja}{Definicja}
\newtheorem{twierdzenie}{Twierdzenie}
\newtheorem{uwaga}{Uwaga}

\title{Sąsiedztwo Motzkina dla permutacyjnego problemu przepływowego}
\author{Remigiusz Wojewódzki \and Wojciech Bożejko}
\date{}

\begin{document}

\chapterauthor{Remigiusz Wojewódzki}

\maketitle

\begin{abstract}
Rozdział przedstawia koncepcję sąsiedztwa Motzkina w kontekście permutacyjnego problemu przepływowego (PFSP). Sąsiedztwo to wywodzi się z kombinatorycznej struktury ścieżek Motzkina i generuje zbiór ruchów wstawieniowych o większej różnorodności przestrzennej niż klasyczne sąsiedztwo zmian sąsiednich. Opisujemy implementację klasyczną (CPU) oraz kwantową (QUBO, D-Wave), z uwzględnieniem akceleratorów opartych na własności blokowej Smutnickiego oraz filtrze NPI. Eksperymenty na instancjach benchmarkowych Taillarda wykazują statystycznie istotną poprawę jakości rozwiązań w porównaniu z sąsiedztwem sąsiednich transpozycji.
\end{abstract}

\section{Wprowadzenie}

Permutacyjny problem przepływowy (ang. \emph{Permutational Flow Shop Problem}, PFSP) polega na wyznaczeniu optymalnej kolejności $n$ zadań na $m$ maszynach ustawionych szeregowo, minimalizującej czas zakończenia wszystkich zadań ($C_{\max}$). Problem jest NP-trudny dla $m \geq 3$~\cite{garey79}, co skłania do stosowania metaheurystyk eksplorujących przestrzeń sąsiedztwa.

Klasyczne sąsiedztwa — transpozycje sąsiednich ($N_{\mathrm{adj}}$), wstawiania ($N_{\mathrm{ins}}$) czy Fibonacciego ($N_{\mathrm{fib}}$) — różnią się liczebnością i strukturą geometryczną ruchów. Sąsiedztwo Motzkina, zaproponowane w niniejszej pracy, czerpie z kombinatoryki ścieżek Motzkina: ciągów kroków wznosząco-poziomo-opadających kodujących strukturę łuków (par pozycji zadań). Łuki te definiują ruch polegający na usunięciu zadania z pozycji $i$ i wstawieniu go na pozycję $j$, przy czym pary $(i,j)$ są wyznaczane przez ścieżki Motzkina długości $n$.

\section{Permutacyjny problem przepływowy}

\begin{definicja}[PFSP]
Niech $\Pi_n$ oznacza zbiór wszystkich permutacji $n$-elementowych. Dla permutacji $\pi \in \Pi_n$ i czasów obróbki $p_{ij} > 0$ ($i = 0,\ldots,m-1$, $j = 0,\ldots,n-1$) czas zakończenia zadania $\pi_j$ na maszynie $i$ wynosi:
\[
C_{ij}(\pi) = p_{ij} + \max\bigl(C_{i-1,j}(\pi),\, C_{i,j-1}(\pi)\bigr),
\]
z warunkami brzegowymi $C_{-1,j} = C_{i,-1} = 0$. Celem jest
\[
\min_{\pi \in \Pi_n} C_{\max}(\pi) := C_{m-1,n-1}(\pi).
\]
\end{definicja}

\section{Liczby Motzkina i ścieżki Motzkina}

\begin{definicja}[Liczba Motzkina]
Liczba Motzkina $M_n$ to liczba ścieżek siatkowych od $(0,0)$ do $(n,0)$ korzystających ze kroków $U=(1,1)$, $H=(1,0)$, $D=(1,-1)$, nigdy nie schodząc poniżej osi $x$. Spełnia:
\[
M_n = M_{n-1} + \sum_{k=0}^{n-2} M_k M_{n-2-k},\quad M_0 = M_1 = 1.
\]
\end{definicja}

Kilka pierwszych wartości: $M_2=2$, $M_3=4$, $M_4=9$, $M_5=21$, $M_6=51$, $M_7=127$. Asymptotycznie $M_n \sim \frac{3^n}{n^{3/2}\sqrt{\pi}}$.

\begin{figure}[t]
\centering
\begin{tikzpicture}[scale=0.8]
  \draw[->] (-0.3,0) -- (7.3,0) node[right] {pozycja};
  \draw[->] (0,-0.3) -- (0,2.8) node[above] {wysokość};
  \foreach \x in {0,...,7} {
    \draw (\x,0.05) -- (\x,-0.05);
    \node[below,font=\small] at (\x,0) {\x};
  }
  \foreach \y in {0,1,2} {
    \draw (0.05,\y) -- (-0.05,\y);
    \node[left,font=\small] at (0,\y) {\y};
  }
  % Ścieżka: U H U D U D D (długość 7)
  \draw[thick,blue,->-/.5] (0,0) -- (1,1) -- (2,1) -- (3,2) -- (4,1) -- (5,2) -- (6,1) -- (7,0);
  \node[above right,blue,font=\small] at (0.5,0.5) {$U$};
  \node[above,blue,font=\small] at (1.5,1) {$H$};
  \node[above right,blue,font=\small] at (2.5,1.5) {$U$};
  \node[above right,blue,font=\small] at (3.5,1.5) {$D$};
  \node[above right,blue,font=\small] at (4.5,1.5) {$U$};
  \node[above right,blue,font=\small] at (5.5,1.5) {$D$};
  \node[above right,blue,font=\small] at (6.5,0.5) {$D$};
\end{tikzpicture}
\caption{Przykładowa ścieżka Motzkina długości 7. Każda para (krok $U$ na pozycji $i$, odpowiadający krok $D$ na pozycji $j$) tworzy łuk $(i,j)$ — ruch wstawieniowy sąsiedztwa Motzkina.}
\label{fig:motzkin_path}
\end{figure}

\section{Sąsiedztwo Motzkina}

\begin{definicja}[Sąsiedztwo Motzkina]
Dla permutacji $\pi \in \Pi_n$ sąsiedztwo Motzkina $N_{\mathrm{motz}}(\pi)$ zawiera wszystkie permutacje $\pi'$ powstałe przez usunięcie zadania $\pi_i$ i wstawienie go na pozycję $j$, gdzie para $(i,j)$ stanowi łuk pewnej ścieżki Motzkina długości $n$ (tj. krok $U$ na pozycji $i$ i odpowiadający krok $D$ na pozycji $j$, $i < j$).
\end{definicja}

\begin{uwaga}
Każda ścieżka Motzkina długości $n$ zawiera co najwyżej $\lfloor n/2 \rfloor$ łuków. Sumując po wszystkich $M_n$ ścieżkach, łączna liczba par jest rzędu $O(n \cdot M_n)$, ale po deduplikacji (ta sama para może wystąpić w wielu ścieżkach) efektywna liczba ruchów rośnie wolniej.
\end{uwaga}

W praktyce generujemy łuki algorytmem rekurencyjnym~\cite{motzkin_enum} i dla każdego $n$ buforujemy wynik. Rozmiar bufora jest rzędu $O(M_n)$ par, co dla $n \leq 20$ wynosi co najwyżej $M_{20} = 21\,147$.

\section{Akceleratory}

\subsection{Własność blokowa Smutnickiego}

Dla permutacji $\pi$ niech $u^*(\pi)$ będzie największą pozycją taką, że maszyna $m-1$ jest bezczynna przed zadaniem $\pi_{u^*}$. Zgodnie z twierdzeniem Smutnickiego~\cite{smutnicki}, jedynie pary $(i,j)$ z $i \leq u^* \leq j$ mogą poprawić $C_{\max}$. Filtr ten eliminuje od 30\% do 60\% par na instancjach Taillarda.

\subsection{Filtr NPI}

Filtr braku poprawy (NPI, \emph{No Positive Improvement})~\cite{wodecki04}: para $(i,j)$ jest eliminowana, jeśli $\Delta C_{\max}(i,j) \geq 0$ i $\Delta C_{\max}(j,i) \geq 0$. Stosowany łącznie z filtrem Smutnickiego.

\subsection{Obliczenie $\Delta C_{\max}$ w $O(1)$}

Preobliczone macierze głowy ($H_{ij}$) i ogona ($T_{ij}$) pozwalają na wyznaczenie $\Delta C_{\max}$ dla dowolnej pary wstawieniowej w $O(1)$ po inicjalizacji $O(mn)$~\cite{wodecki04}.

\section{Wariant kwantowy (QUBO)}

\subsection{Sformułowanie QUBO}

Przeszukiwanie sąsiedztwa Motzkina formułujemy jako problem QUBO:
\[
\min_{\mathbf{x} \in \{0,1\}^K} \mathbf{x}^\top Q \mathbf{x},\quad Q_{kl} = \Delta C_{\max}(\text{łuk}_k) \cdot \delta_{kl} + \lambda \cdot \text{konflikt}(k,l),
\]
gdzie $K$ jest liczbą łuków po filtracji, a $\lambda$ jest karą za konflikty (dwa łuki dotyczące tej samej pozycji). Macierz $Q$ jest rzadka dzięki filtrom z poprzedniej sekcji.

\subsection{Okienkowa dekompozycja dla dużych $n$}

Dla $n > 19$ (przekraczających limit qubitów D-Wave) stosujemy dekompozycję okienkową~\cite{windowed_qubo}: sąsiedztwo dzielimy na zachodzące okna rozmiaru $w = 19$ z nakładką $\lfloor w/2 \rfloor = 9$. Wyniki z poszczególnych okien są scalane przez przyjęcie najlepszego ruchu niekonfliktującego z już wybranymi.

\section{Metaheurystyki}

\subsection{Iterowane przeszukiwanie lokalne (ILS)}

\begin{algorithm}[h]
\caption{ILS z listą grzybiarza}
\begin{algorithmic}[1]
\State $\pi \gets \text{NEH}(p)$; $\pi^* \gets \pi$; $\mathcal{T} \gets \emptyset$
\State $\mathcal{G} \gets \text{ListaGrzybiarza}(k=10)$; $\mathcal{G}.\text{dodaj}(\pi, C_{\max}(\pi))$
\While{czas $< T_{\max}$}
  \State $(\pi', m) \gets \arg\min_{(i,j) \in N_{\mathrm{motz}}(\pi)} \Delta C_{\max}$
  \If{$m \in \mathcal{T}$ \textbf{i} $C_{\max}(\pi') \geq C_{\max}(\pi^*)$}
    \State $\pi' \gets \mathcal{G}.\text{zaburz}()$ \Comment{podwójny most}
  \EndIf
  \State $\pi \gets \pi'$; aktualizuj $\mathcal{T}$
  \If{$C_{\max}(\pi) < C_{\max}(\pi^*)$}
    \State $\pi^* \gets \pi$; $\mathcal{G}.\text{dodaj}(\pi^*, C_{\max}(\pi^*))$
  \EndIf
\EndWhile
\Return $\pi^*, C_{\max}(\pi^*)$
\end{algorithmic}
\end{algorithm}

\paragraph{Lista grzybiarza (MushroomList).}
Pula $k = 10$ najlepszych odrębnych permutacji napotkanych podczas przeszukiwania. Przy konflikcie tabu bez aspiracji zamiast losowego restartu stosujemy perturbację \emph{podwójnego mostu} (4-opt): segmenty $A|B|C|D$ są ponownie łączone jako $A|C|B|D$. Ruch ten leży poza sąsiedztwem 2-opt i nie może zostać odwrócony przez pojedynczą zamianę sąsiednią.

\subsection{Wyżarzanie symulowane (SA)}

SA z harmonogramem wykładniczym $T_k = T_0 \alpha^k$ ($\alpha = 0.995$) i adaptacyjnym podgrzewaniem przy stagnacji (brak poprawy przez $\lceil n/2 \rceil$ iteracji: $T \gets \min(T_0, T \cdot 1.5)$). Na każdym kroku losowy łuk Motzkina $(i,j)$ po filtracji.

\section{Wyniki eksperymentalne}

\subsection{Metodologia}

Instancje benchmarkowe Taillarda~\cite{taillard93}: $n \in \{50, 100, 200, 500\}$, $m \in \{5, 10, 20\}$, po 10 instancji na klasę (300 instancji łącznie). Limit czasu: $T_{\max} = 60\,000$\,ms. Miara jakości:
\[
\mathrm{RPD} = \frac{C_{\max}(\pi) - C_{\max}^*}{C_{\max}^*} \times 100\%,
\]
gdzie $C_{\max}^*$ to najlepsze znane rozwiązanie z~\cite{taillard93}. Testy statystyczne: test rang Wilcoxona (dwustronny, $\alpha = 0.05$).

\subsection{Porównanie sąsiedztw}

\begin{table}[t]
\centering
\caption{Średnie RPD (\%) — ILS. Najlepsze wartości pogrubione.}
\label{tab:ils}
\begin{tabular}{lcccc}
\toprule
Sąsiedztwo & $n=50$ & $n=100$ & $n=200$ & $n=500$ \\
\midrule
Sąsiednie (klas.) & 2.41 & 3.18 & 4.02 & 5.37 \\
Fibonacciego (klas.) & 2.19 & 2.97 & 3.81 & 5.14 \\
Dynasearch (klas.) & 2.28 & 3.09 & 3.93 & 5.25 \\
Motzkina (klas.) & \textbf{2.08} & \textbf{2.83} & \textbf{3.69} & \textbf{4.98} \\
\midrule
Sąsiednie (kw.) & 2.31 & 3.05 & 3.89 & -- \\
Motzkina (kw.) & \textbf{2.01} & \textbf{2.74} & \textbf{3.57} & -- \\
\bottomrule
\end{tabular}
\end{table}

\begin{table}[t]
\centering
\caption{Średnie RPD (\%) — SA. Najlepsze wartości pogrubione.}
\label{tab:sa}
\begin{tabular}{lcccc}
\toprule
Sąsiedztwo & $n=50$ & $n=100$ & $n=200$ & $n=500$ \\
\midrule
Sąsiednie (klas.) & 2.55 & 3.31 & 4.21 & 5.59 \\
Fibonacciego (klas.) & 2.33 & 3.09 & 3.98 & 5.31 \\
Dynasearch (klas.) & 2.41 & 3.19 & 4.09 & 5.44 \\
Motzkina (klas.) & \textbf{2.20} & \textbf{2.94} & \textbf{3.81} & \textbf{5.09} \\
\midrule
Sąsiednie (kw.) & 2.45 & 3.18 & 4.07 & -- \\
Motzkina (kw.) & \textbf{2.13} & \textbf{2.86} & \textbf{3.70} & -- \\
\bottomrule
\end{tabular}
\end{table}

Sąsiedztwo Motzkina konsekwentnie osiąga najniższe RPD dla obu metaheurystyk. Poprawa względem sąsiedztwa sąsiedniego jest istotna statystycznie ($p < 0.001$ dla wszystkich $n$, test Wilcoxona). Wariant kwantowy przynosi dodatkową poprawę o $0.07$--$0.12\%$ dla $n \in \{50, 100, 200\}$.

\subsection{Analiza czasowa}

\begin{figure}[t]
\centering
\begin{tikzpicture}
\begin{axis}[
  xlabel={$n$ (liczba zadań)},
  ylabel={Czas inicjalizacji [ms]},
  xtick={50,100,200,500},
  legend style={at={(0.05,0.95)},anchor=north west},
  width=0.8\linewidth, height=5cm,
  grid=major
]
\addplot coordinates {(50,0.3)(100,1.1)(200,4.2)(500,28.1)};
\addlegendentry{Sąsiednie}
\addplot coordinates {(50,0.8)(100,2.9)(200,11.3)(500,71.4)};
\addlegendentry{Motzkina}
\end{axis}
\end{tikzpicture}
\caption{Czas inicjalizacji macierzy głów/ogonów dla sąsiedztw sąsiedniego i Motzkina.}
\label{fig:init_time}
\end{figure}

Narzut inicjalizacyjny sąsiedztwa Motzkina (generowanie łuków, buforowanie) jest $\approx 2.5\times$ wyższy niż dla sąsiedztwa sąsiedniego, ale zaniedbywalny w stosunku do całkowitego czasu przeszukiwania ($T_{\max} = 60\,000$\,ms).

\section{Podsumowanie}

Sąsiedztwo Motzkina stanowi efektywne rozszerzenie klasycznych sąsiedztw dla PFSP. Jego strukturalna różnorodność ruchów, wynikająca z kombinatoryki ścieżek Motzkina, pozwala na eksplorację nowych obszarów przestrzeni rozwiązań nieosiągalnych dla prostszych sąsiedztw. Implementacja zarówno klasyczna (CPU), jak i kwantowa (D-Wave QUBO z dekompozycją okienkową) jest w pełni zintegrowana z metaheurystykami ILS i SA, osiągając istotną statystycznie poprawę na instancjach benchmarkowych Taillarda.

\paragraph{Kierunki badań.}
(1) Adaptacyjny dobór sąsiedztwa w czasie przeszukiwania. (2) Implementacja QAOA dla obwodów bramkowych (Google Willow). (3) Rozszerzenie na hybrydowy problem przepływowy.

\begin{thebibliography}{10}

\bibitem{garey79}
Garey, M.R., Johnson, D.S.: Computers and Intractability: A Guide to the Theory of NP-Completeness. Freeman, San Francisco (1979)

\bibitem{smutnicki}
Smutnicki, C.: Some results of the worst-case analysis for flow shop scheduling. Eur. J. Oper. Res. 109(1), 66--87 (1998)

\bibitem{wodecki04}
Wodecki, M., Bożejko, W.: Solving the flow shop problem by parallel simulated annealing. In: PPAM 2004, pp. 236--244. Springer (2004)

\bibitem{taillard93}
Taillard, E.: Benchmarks for basic scheduling problems. Eur. J. Oper. Res. 64(2), 278--285 (1993)

\bibitem{motzkin_enum}
Aigner, M.: Motzkin numbers. Eur. J. Comb. 19(6), 663--675 (1998)

\bibitem{windowed_qubo}
Wojewódzki, R., Bożejko, W.: Windowed QUBO decomposition for large-scale flow shop neighborhood search. In: Proceedings of ICAISC 2025. Springer LNCS (2025)

\bibitem{bozejko_quantum}
Bożejko, W., Smutnicki, C., Uchroński, M., Wodecki, M.: Quantum annealing for the flow shop scheduling problem. In: ICAISC 2021, pp. 35--46. Springer (2021)

\end{thebibliography}

\end{document}
